
\hspace*{0.3in} The rapid proliferation of cloud-native applications has fundamentally transformed how modern software systems are architected and deployed. Organizations increasingly rely on microservice-based architectures — where a single application is decomposed into tens or hundreds of independently deployable services communicating over network boundaries — to achieve scalability, fault isolation, and continuous delivery. While this paradigm offers compelling engineering advantages, it simultaneously introduces an extraordinary degree of operational complexity. When a failure occurs in such an environment, the root cause rarely manifests at the point of observation; instead, cascading dependency chains propagate anomalies across service boundaries, making the failure surface appear diffuse and deceptive. Site Reliability Engineers (SREs) tasked with diagnosing such incidents often spend hours correlating logs, metrics, and traces across dozens of services, a process that is both time-consuming and heavily dependent on institutional knowledge that is difficult to retain and transfer.\\\\

\hspace*{0.3in} SYNAPSE — AI-Powered AIOps Platform for Automated Root Cause Analysis of Microservice Failures — is designed to address this operational challenge at its core. The platform continuously ingests real-time telemetry from all services in a deployment, constructs dynamic service dependency graphs, and applies a hybrid deep learning pipeline to identify, localize, and explain the root causes of incidents without requiring manual intervention. At the heart of SYNAPSE lies a multi-stage AI pipeline: a Variational Autoencoder coupled with a Graph Neural Network (VAE-GNN) performs topology-aware anomaly detection by modeling the normal joint distribution of metrics across the dependency graph; a Differentiable Causal Discovery module (DECI), constrained by Retrieval-Augmented Generation (RAG) priors derived from historical incident knowledge, infers the causal structure underlying each anomaly; and a Large Language Model (LLM) reasoning layer synthesizes all upstream signals into human-readable diagnostic narratives and ranked remediation recommendations. The result is a closed-loop AIOps system capable of transforming raw observability data into actionable root cause conclusions within seconds of anomaly onset.\\\\

\hspace*{0.3in} The practical impact of SYNAPSE extends beyond raw speed of diagnosis. By encoding causal relationships rather than mere correlations, the platform dramatically reduces false-positive alert fatigue and correctly attributes cascading failures to their originating service even when observable symptoms are spatially and temporally displaced from the true cause. Automated recovery action selection, implemented through a curated kubectl-based remediation catalog, allows the system to trigger targeted corrective actions — pod restarts, traffic rerouting, resource quota adjustments — immediately upon high-confidence root cause identification. Critically, SYNAPSE incorporates a continual learning subsystem based on Elastic Weight Consolidation (EWC) and Model-Agnostic Meta-Learning (MAML), enabling the platform to assimilate feedback from newly resolved incidents and adapt to evolving service topologies without catastrophic forgetting of previously learned failure patterns. Taken together, these capabilities position SYNAPSE as a comprehensive, production-ready AIOps platform that materially reduces Mean Time to Resolution (MTTR) and raises the reliability ceiling of microservice deployments.\\\\

\section{Project Idea}

\hspace*{0.3in} The central idea behind SYNAPSE is to construct an end-to-end, AI-driven pipeline that mirrors the diagnostic reasoning process of an experienced SRE, but operates at machine speed and scale. The system begins with a high-frequency telemetry ingestion engine that collects CPU utilization, memory consumption, request latency, error rates, saturation metrics, and inter-service call volumes from all microservices in the target deployment. These raw time-series streams are not treated independently; instead, they are organized into a service dependency graph whose topology is either discovered automatically through distributed tracing instrumentation or supplied via configuration. This graph serves as the structural backbone for every subsequent reasoning step, ensuring that spatial relationships between services are preserved and exploited throughout the analysis.\\\\

\hspace*{0.3in} Anomaly detection in SYNAPSE is performed by a VAE-GNN model that jointly encodes the graph topology and the multivariate metric time-series into a shared latent space. The encoder learns the normal operational manifold of the entire service mesh during a training phase on historical telemetry; at inference time, deviations from this manifold — quantified as reconstruction error and latent-space displacement — are used to flag anomalous services and assign them a severity level on a four-tier scale: L1 (Simple), L2 (Cascading), L3 (Systemic), and L4 (Design Flaw). Once anomalies are detected, SYNAPSE activates its causal discovery engine, DECI, which frames the problem of identifying the root cause as a Bayesian structure learning task over the service dependency graph. DECI's optimization is guided by a RAG-constructed Weighted Prior ($W_0$) that encodes domain knowledge from semantically similar historical incidents retrieved from the platform's pattern history database, preventing the causal search from converging to implausible structures and substantially reducing the solution space.\\\\

\hspace*{0.3in} The root cause and causal graph produced by DECI are then passed to an LLM reasoning module powered by GPT-4o-mini. The LLM receives a structured prompt containing the anomaly profile, the inferred causal chain, relevant historical context retrieved via RAG, and the current severity classification, and generates a coherent diagnostic explanation alongside a ranked list of recovery actions. These actions are mapped to concrete kubectl commands stored in the recovery catalog and can be executed automatically upon operator approval or in fully autonomous mode depending on deployment policy. The four key technical techniques that underpin this pipeline are:\\\\

\begin{itemize}
    \item \textbf{VAE-GNN Anomaly Detection:} A hybrid deep learning model that fuses Variational Autoencoder-based metric reconstruction with Graph Neural Network message passing to perform topology-aware, joint anomaly detection across the entire service dependency graph.
    \item \textbf{DECI Causal Inference with RAG-Constrained Prior:} A differentiable causal structure learning algorithm that discovers directed acyclic causal graphs over service metrics, whose search is constrained by a weighted adjacency prior constructed from semantically retrieved historical incident patterns, improving both accuracy and convergence speed.
    \item \textbf{LLM Reasoning with Retrieval-Augmented Generation:} A GPT-4o-mini-powered reasoning layer that synthesizes anomaly evidence, causal structure, and retrieved incident knowledge into natural-language root cause explanations and ranked, actionable recovery recommendations.
    \item \textbf{Continual Learning via EWC and MAML:} A dual-strategy continual learning subsystem that uses Elastic Weight Consolidation to prevent catastrophic forgetting of prior failure patterns while applying Model-Agnostic Meta-Learning to enable rapid adaptation to new incident types with minimal retraining data.
\end{itemize}

\hspace*{0.3in} The design philosophy of SYNAPSE is grounded in three overarching principles: causality over correlation, interpretability over opacity, and adaptability over rigidity. Unlike traditional monitoring solutions that rely on threshold-based alerting and statistical correlation to identify related events, SYNAPSE explicitly models the causal mechanisms through which one service failure propagates to and manifests as observable symptoms in downstream services. This causal grounding is essential for correctly attributing responsibility in complex cascading failure scenarios where dozens of services may simultaneously exhibit anomalous behavior, and naive correlation approaches would implicate all of them equally.\\\\

\hspace*{0.3in} Equally central to SYNAPSE's design is the commitment to transparency and human-AI collaboration. Every root cause conclusion produced by the platform is accompanied by a traceable causal graph, a plain-language narrative explaining the chain of events, and confidence scores for each step of the reasoning process. This transparency allows SRE teams to audit, validate, and override the system's conclusions, fostering trust and enabling the platform to serve as a collaborative diagnostic partner rather than an inscrutable black box. The continual learning architecture further ensures that this collaboration is bidirectional: corrections and feedback provided by operators are systematically incorporated into the model's knowledge, making SYNAPSE progressively more accurate and aligned with the specific operational characteristics of each deployment it monitors.\\\\

\section{Motivation of the Project}

\hspace*{0.3in} The motivation for building SYNAPSE emerges from a convergence of practical operational pain points observed in real-world microservice deployments and fundamental limitations in the current generation of monitoring and observability tools. As organizations scale their microservice estates into hundreds of independently deployed components, the gap between the volume and velocity of telemetry data generated and the human capacity to interpret it grows into an unbridgeable chasm. The following specific challenges served as the primary drivers for the conception and development of the SYNAPSE platform:\\\\

\begin{enumerate}
    \item \textbf{Manual Root Cause Analysis is Prohibitively Slow:} In production microservice environments, the average Mean Time to Resolution (MTTR) for incidents caused by service-level failures routinely exceeds 30–60 minutes, with complex cascading failures sometimes remaining unresolved for hours. SREs must manually correlate metrics, logs, and distributed traces across dozens of dashboards, a process that is inherently sequential and bottlenecked by human cognitive bandwidth, creating an urgent need for automated diagnostic acceleration.

    \item \textbf{Existing Tools Cannot Localize Cascading Failures Accurately:} Current AIOps and observability platforms predominantly rely on metric correlation and static alerting rules to surface incidents. These approaches systematically fail in cascading failure scenarios — where a single root-cause event in one service triggers a chain of symptomatic anomalies in dependent services — because they attribute equal diagnostic weight to every anomalous service rather than tracing the failure back to its point of origin.

    \item \textbf{Lack of Automated Causal Inference in Operational Tooling:} The distinction between correlation and causation is critical in failure diagnosis, yet virtually no production-grade observability tool incorporates genuine causal discovery. Without a causal model of service interactions, operators cannot reliably distinguish a service that is the origin of a failure from one that is merely a downstream victim exhibiting sympathetic symptoms, leading to misattribution and wasted remediation effort.

    \item \textbf{Diagnostic Outputs Lack Human Interpretability:} Even when machine learning models are applied to anomaly detection, their outputs are typically limited to anomaly scores or binary flags that provide no contextual explanation of why a particular service is flagged or what sequence of events led to the current state. On-call engineers require narrative explanations that map technical evidence to operational understanding, a capability that demands the integration of natural language generation into the diagnostic pipeline.

    \item \textbf{Recovery Actions Remain Manual and Inconsistent:} After root cause identification, the remediation process in most organizations is entirely manual, relying on runbooks that may be outdated, incomplete, or unknown to the responding engineer. This introduces further delays, increases the risk of incorrect remediation, and makes incident response quality dependent on the specific expertise of whoever is on call rather than on systematically accumulated organizational knowledge.

    \item \textbf{Models Degrade as Service Topologies Evolve:} Static machine learning models trained on historical telemetry quickly become stale as microservice deployments evolve — new services are added, dependencies change, and traffic patterns shift with product releases. Without a mechanism for continual learning, anomaly detection and causal inference models accumulate drift over time, producing increasing rates of false positives and missed detections that erode operator trust in the system.

    \item \textbf{Absence of a Unified, End-to-End AIOps Platform:} Organizations attempting to address the above challenges are forced to integrate and maintain a fragmented collection of disparate tools — separate systems for metric collection, anomaly detection, log analysis, tracing, and runbook automation — that were not designed to work together and whose outputs cannot be coherently synthesized into a unified diagnostic conclusion. The operational overhead of maintaining this toolchain, coupled with the absence of a shared data model across tools, constitutes a significant engineering burden that diverts resources away from reliability improvements.
\end{enumerate}

\hspace*{0.3in} SYNAPSE was conceived as a direct and comprehensive response to all seven of the above challenges within a single, cohesive platform. By unifying telemetry ingestion, graph-structured anomaly detection, causal discovery, LLM-driven explanation generation, automated recovery action selection, and continual learning into a coherent, tightly integrated pipeline, SYNAPSE eliminates the integration complexity of fragmented toolchains while simultaneously delivering diagnostic capabilities that no individual existing tool provides. The platform is designed to be deployment-agnostic and incrementally adoptable, allowing organizations to integrate it into existing Kubernetes-based environments with minimal instrumentation overhead and immediately begin realizing reductions in MTTR and improvements in incident attribution accuracy.\\\\
